\documentclass[aspectratio=169,11pt]{beamer}
% Shared CMPSC 480 Beamer preamble.
\usepackage[T1]{fontenc}
\usepackage[utf8]{inputenc}
\usepackage{lmodern}
\usepackage{amsmath,amssymb,mathtools}
\usepackage{booktabs,array,tabularx}
\usepackage{xcolor}
\usepackage{listings}
\usepackage{tikz}
\usetikzlibrary{arrows.meta,positioning,shapes.geometric}

% Ignore only negligible TeX box diagnostics that are below visual tolerance.
% Larger layout defects still appear in the build log and in rendered QA.
\vfuzz=3pt
\hbadness=2000

\definecolor{courseink}{HTML}{111111}
\definecolor{coursegray}{HTML}{EDEDED}
\definecolor{courserule}{HTML}{B8BCC4}
\definecolor{courseblue}{HTML}{3D8DFF}
\definecolor{courselightblue}{HTML}{D0EDFA}
\definecolor{coursemuted}{HTML}{5C6470}
\definecolor{coursegreen}{HTML}{0B7A53}
\definecolor{coursered}{HTML}{B3261E}

\usetheme{default}
\usefonttheme{professionalfonts}
\setbeamertemplate{navigation symbols}{}
\setbeamersize{text margin left=0.65cm,text margin right=0.65cm}

\setbeamercolor{normal text}{fg=courseink,bg=white}
\setbeamercolor{frametitle}{fg=courseink,bg=white}
\setbeamercolor{title}{fg=courseink,bg=white}
\setbeamercolor{subtitle}{fg=coursemuted,bg=white}
\setbeamercolor{structure}{fg=courseblue}
\setbeamercolor{alerted text}{fg=coursered}
\setbeamercolor{block title}{fg=courseink,bg=courselightblue}
\setbeamercolor{block body}{fg=courseink,bg=coursegray}
\setbeamercolor{example text}{fg=coursegreen}

\setbeamerfont{title}{size=\fontsize{25}{29}\selectfont,series=\bfseries}
\setbeamerfont{subtitle}{size=\fontsize{13}{16}\selectfont}
\setbeamerfont{frametitle}{size=\fontsize{19}{22}\selectfont,series=\bfseries}
\setbeamerfont{normal text}{size=\fontsize{11}{14}\selectfont}
\setbeamerfont{footline}{size=\fontsize{7.5}{9}\selectfont}

\setbeamertemplate{frametitle}{%
  \nointerlineskip
  % Use content-driven height so a long assessment title can wrap without
  % being clipped by a fixed-height title bar.
  \begin{beamercolorbox}[wd=\paperwidth,sep=0.17cm,leftskip=0.48cm,rightskip=0.48cm]{frametitle}
    \insertframetitle
  \end{beamercolorbox}
  \vspace{-0.08cm}
  {\color{courserule}\rule{\textwidth}{0.35pt}}
}

\setbeamertemplate{footline}{%
  \leavevmode
  \begin{beamercolorbox}[wd=\paperwidth,ht=2.6ex,dp=1.1ex,leftskip=.65cm,rightskip=.65cm]{author in head/foot}
    \color{coursemuted}\insertshorttitle\hfill\insertframenumber/\inserttotalframenumber
  \end{beamercolorbox}
}

\setbeamertemplate{itemize item}{\color{courseblue}\small$\blacktriangleright$}
\setbeamertemplate{itemize subitem}{\color{courseblue}\scriptsize$\bullet$}
\setbeamertemplate{enumerate item}{\color{courseblue}\insertenumlabel.}

\lstdefinestyle{coursepython}{
  language=Python,
  basicstyle=\ttfamily\fontsize{8.5}{10}\selectfont,
  keywordstyle=\color{courseblue}\bfseries,
  commentstyle=\color{coursemuted}\itshape,
  stringstyle=\color{coursegreen},
  showstringspaces=false,
  columns=fullflexible,
  keepspaces=true,
  frame=single,
  rulecolor=\color{courserule},
  backgroundcolor=\color{coursegray},
  breaklines=true,
  tabsize=4,
  xleftmargin=0.15cm,
  xrightmargin=0.15cm,
  framexleftmargin=0.15cm,
  framexrightmargin=0.15cm
}
\lstset{style=coursepython}

\newcommand{\points}[1]{\hfill{\color{courseblue}\textbf{[#1 points]}}}
\newcommand{\deliverable}[1]{\textbf{\color{courseblue}#1}}
\newcommand{\code}[1]{\texttt{#1}}
\newcommand{\keyidea}[1]{%
  \begin{beamercolorbox}[rounded=false,sep=0.55em,wd=\linewidth]{block body}
    \textbf{Key idea.} #1
  \end{beamercolorbox}
}
\newcommand{\answerlabel}{\textbf{\color{coursegreen}Answer.}\ }
\newcommand{\gradinglabel}{\textbf{\color{courseblue}Grading guidance.}\ }

\newenvironment{questionblock}[1]
  {\begin{block}{#1}}
  {\end{block}}

\newenvironment{solutionblock}[1]
  {\begin{block}{#1}}
  {\end{block}}

\AtBeginSection[]{
  \begin{frame}
    \vfill
    {\usebeamerfont{title}\color{courseink}\insertsectionhead\par}
    \vspace{0.4cm}
    {\color{courseblue}\rule{0.32\paperwidth}{3pt}}
    \vfill
  \end{frame}
}


% CMPSC480_TOTAL_POINTS: 20
% CMPSC480_ITEM_IDS: 1,2,3,4
\title[Quiz 2 Solutions]{Quiz 2: Adversarial and Sequential Decision Making}
\subtitle{Weeks 3-5}
\author{CMPSC 480 - Instructor Solution Deck}
\date{}

\begin{document}


\begin{frame}
  \titlepage
\vspace{-0.35cm}
\begin{center}
\color{coursered}\bfseries Instructor material - do not distribute with the question deck
\end{center}
\end{frame}

\begin{frame}{Solution overview}
\begin{columns}[T,onlytextwidth]
\begin{column}{0.62\textwidth}
\Large This instructor deck provides a complete matched solution and grading guidance.

\vspace{0.35cm}
\normalsize
Coverage: Weeks 3-5
\end{column}
\begin{column}{0.33\textwidth}
\begin{block}{Points}20 total\end{block}
\begin{block}{Expected time}35 minutes\end{block}
\begin{block}{Submission}Individual PDF\end{block}
\end{column}
\end{columns}
\end{frame}

\begin{frame}{Instructor verification checklist}
\small
\begin{itemize}
  \item Work independently.
  \item Show backed-up values and arithmetic.
  \item State terminal conventions.
  \item No electronic communication.
\end{itemize}
\end{frame}

\begin{frame}{Solution 1.a - Minimax and alpha-beta (2 points)}
\small
\textbf{Question focus.} Compute the minimax value and selected branch.

\vspace{0.25cm}
\answerlabel MIN(A)=3 and MIN(B)=2, so MAX chooses A with value 3.
\end{frame}

\begin{frame}{Solution 1.b - Minimax and alpha-beta (2 points)}
\small
\textbf{Question focus.} Search A first. After B's first leaf 2, can B's second leaf be pruned?

\vspace{0.25cm}
\answerlabel Yes. The root alpha is 3 after A; B's beta becomes 2, so alpha is at least beta and the remaining B child cannot improve MAX's decision.
\end{frame}

\begin{frame}{Solution 1.c - Minimax and alpha-beta (2 points)}
\small
\textbf{Question focus.} Does pruning change the backed-up value? Explain.

\vspace{0.25cm}
\answerlabel No. Alpha-beta removes branches proven irrelevant to the root decision and returns the same minimax value under exact arithmetic.
\end{frame}

\begin{frame}{Solution 2.a - MDP backup (2 points)}
\small
\textbf{Question focus.} Write the Bellman action-value expression.

\vspace{0.25cm}
\answerlabel Q(s,a)=0.8[0+0.9*5]+0.2[0+0.9*2] under the convention that terminal state's listed value is 5.
\end{frame}

\begin{frame}{Solution 2.b - MDP backup (2 points)}
\small
\textbf{Question focus.} Compute Q(s,a).

\vspace{0.25cm}
\answerlabel The result is 0.8*4.5+0.2*1.8=3.6+0.36=3.96.
\end{frame}

\begin{frame}{Solution 2.c - MDP backup (2 points)}
\small
\textbf{Question focus.} Explain what value iteration maximizes.

\vspace{0.25cm}
\answerlabel For every nonterminal state it computes this expected return for each legal action and sets the new value to the maximum.
\end{frame}

\begin{frame}{Solution 3.a - Q-learning (2 points)}
\small
\textbf{Question focus.} Compute the nonterminal target and TD error.

\vspace{0.25cm}
\answerlabel Target=2+0.5*3=3.5 and delta=3.5-1=2.5.
\end{frame}

\begin{frame}{Solution 3.b - Q-learning (2 points)}
\small
\textbf{Question focus.} Compute the updated Q value and state the terminal change.

\vspace{0.25cm}
\answerlabel New Q=1+0.25*2.5=1.625. If terminal, target is reward 2 and new Q is 1.25.
\end{frame}

\begin{frame}{Solution 4.a - Planning versus learning (2 points)}
\small
\textbf{Question focus.} What information does each method require?

\vspace{0.25cm}
\answerlabel Value iteration requires an explicit transition and reward model. Q-learning can learn from sampled transitions without knowing those probabilities.
\end{frame}

\begin{frame}{Solution 4.b - Planning versus learning (2 points)}
\small
\textbf{Question focus.} Why must a Q-learning evaluation run use epsilon zero?

\vspace{0.25cm}
\answerlabel Exploration intentionally chooses suboptimal actions, so leaving epsilon positive confounds learned policy quality with exploration noise.
\end{frame}

\begin{frame}{Edge-case reasoning and expected behavior}
\small
\begin{itemize}
  \item Return utility without expanding.
  \item Reject or normalize only under an explicit data-cleaning policy.
  \item Remove the bootstrap term.
\end{itemize}
\end{frame}

\begin{frame}{Grading rubric}
\scriptsize
\begin{tabularx}{\textwidth}{>{\bfseries}p{0.28\textwidth} p{0.08\textwidth} X}
\toprule
Category & Points & Full-credit evidence \\
\midrule
Minimax and alpha-beta & 6 & Backups and pruning reasoning are correct. \\
MDP backup & 6 & Probability, reward, and discount are combined correctly. \\
Q-learning & 4 & Target, TD error, and update are correct. \\
Planning versus learning & 4 & Model, data, and convergence distinctions are precise. \\
\bottomrule
\end{tabularx}
\end{frame}

\begin{frame}{Reflection exemplar}
\small
\answerlabel Valid reasons include finite exploration, noisy estimates, hyperparameter effects, or a mismatch in transition, reward, or terminal semantics.

\vspace{0.25cm}
\gradinglabel Award credit for a specific claim supported by technical evidence from the submitted work.
\end{frame}

\end{document}
